%%%PAGE 135 rot=0 legibility=fair summary=双星结论(质心杠杆平衡、周期公式、等效质量)；三星模型：共线三星、正三角形三星、3m+m+m 三星的余弦定理与质心法
%%%BLOCK id=135-01 ch=05 sec=双星模型 kind=formula alt=none cont=none y=0.03-0.30
\textbf{双星结论.}\quad 双星及三星
\[
T=2\pi\sqrt{\frac{L^{3}_{\text{距离}}}{Gm_{\text{和}}}}\qquad
\omega=\sqrt{\frac{GM_{\text{和}}}{L_{\text{距离}}^{3}}}
\] % ω 式分子的质量符号形似大写 M
\[ m_1r_1=m_2r_2 \]

\textbf{一、双星}
%FIG p135_a 0.12 0.155 0.41 0.235
\notefig[0.3]{figures/p135_a.png}
\[
\left\{\begin{aligned}
\frac{Gm_1m_2}{(r_1+r_2)^{2}}&=\frac{m_{1}4\pi^{2}}{T^{2}}r_1\\[4pt]
\frac{Gm_1m_2}{(r_1+r_2)^{2}}&=\frac{m_{2}4\pi^{2}}{T^{2}}r_2
\end{aligned}\right.
\ \longrightarrow\
\left\{\begin{aligned}
&\text{质心\ 杆杠平衡}\\
&m_1r_1=m_2r_2\\[4pt]
&T=\sqrt{\frac{4\pi^{2}(r_1+r_2)^{3}}{G(m_1+m_2)}}
\end{aligned}\right.
\] % 原稿字序为“杆杠平衡”（通常说“杠杆”）；“衡”字后另有一个形如“]”的笔画

\[ m_1\gg m_2\qquad r_1\ll r_2\qquad m_1+m_2=m_1\ \to\ \text{中心天体\unclear{型}} \]
%%%END
%%%BLOCK id=135-02 ch=05 sec=双星模型·等效质量 kind=derivation alt=none cont=none y=0.30-0.47
\textbf{等效质量}
%FIG p135_b 0.05 0.295 0.31 0.354
\notefig[0.3]{figures/p135_b.png}
周期为$T$
%FIG p135_c 0.36 0.29 0.65 0.365
\notefig[0.35]{figures/p135_c.png}
\[
\left\{\begin{array}{l}
\dfrac{Gm_1m_2}{(r_1+r_2)^{2}}=\dfrac{Gm_1\text{\unclear{$M$}}}{r_1^{2}}\\[10pt]
m_1r_1=m_2r_2
\end{array}\right.
\qquad \frac{m_2}{m_1+m_2}=\frac{r_1}{r_1+r_2}
\] % 右式分子的 M 原稿形似小写 m
\[ \therefore\ M=\left(\frac{r_1}{r_1+r_2}\right)^{3}m_2=\frac{m_2^{3}}{(m_1+m_2)^{2}} \] % CHECK: 括号外指数原稿为 3，按推导（与末式 m_2^3/(m_1+m_2)^2 一致）应为 2
等效质量
%%%END
%%%BLOCK id=135-03 ch=05 sec=三星模型·共线三星 kind=derivation alt=none cont=none y=0.47-0.58
\textbf{二、三星}

(1)
%FIG p135_d 0.15 0.484 0.455 0.53
\notefig[0.3]{figures/p135_d.png}
\[
\begin{aligned}
\text{对}m_1:&\quad \frac{Gm_1m_2}{(r_1+r_2)^{2}}+\frac{Gm_1m_0}{r_1^{2}}=\frac{m_{1}4\pi^{2}r_1}{T^{2}}\\[6pt]
\text{对}m_2:&\quad \frac{Gm_1m_2}{(r_1+r_2)^{2}}+\frac{Gm_2m_0}{r_2^{2}}=\frac{m_{2}4\pi^{2}r_2}{T^{2}}
\end{aligned}
\] % 原稿两式共用左侧一个大括号；第二式右端 r_2 右上有一小撇（疑为笔迹）
%%%END
%%%BLOCK id=135-04 ch=05 sec=三星模型·正三角形 kind=derivation alt=none cont=none y=0.58-0.76
(2)\quad 找对称\unclear{类}
%FIG p135_e 0.095 0.595 0.36 0.722
\notefig[0.3]{figures/p135_e.png}
\[ \sqrt{3}\,\frac{Gmm}{a^{2}}=\frac{m4\pi^{2}}{T^{2}}\cdot\frac{\sqrt{3}}{3}\,a \]
\[ T=\sqrt{\frac{4\pi^{2}a^{3}}{3GM}} \] % CHECK: 分母质量符号原稿形似大写 M（上式为 m）
%%%END
%%%BLOCK id=135-05 ch=05 sec=三星模型·余弦定理与质心 kind=derivation alt=none cont=none y=0.76-1.0
(3)\quad 余弦定理法
%FIG p135_f 0.08 0.772 0.425 0.935
\notefig[0.35]{figures/p135_f.png}
\[ r=\sqrt{\left(\frac{a}{2}\right)^{2}+\left(\frac{3}{5}\sqrt{3}\,a\right)^{2}}=\frac{\sqrt{13}}{\sqrt{25}}\,a \] % CHECK: 第二项括号内疑漏 1/2 因子（几何上应为 (3/5)(√3/2)a）

质心位：两两一组找质心\qquad 三星周期相等

令 $\dfrac{Gmm}{a^{2}}=F$ % 此式右侧另有一条斜线(分隔符)
%FIG p135_g 0.715 0.836 0.825 0.905
\notefig[0.2]{figures/p135_g.png}
\red{余弦Th算合力}
\[ F_{\text{合}}=\sqrt{(3F)^{2}+F^{2}\,\text{\unclear{$-$}}\,2F\cdot3F\cos120^\circ}=\sqrt{13}\,F \] % 根号内 F^2 与 2F·3F 之间的符号被红字覆盖（按余弦定理应为 −）
\red{由 $F_{\text{合}}=\dfrac{m4\pi^{2}}{T^{2}}\,\text{\unclear{$r$}}$} % 红字压在上一式下方，末尾的 r 不清
\[ \sqrt{13}\,\frac{Gmm}{a^{2}}=\frac{m4\pi^{2}}{T^{2}}\underbrace{\sqrt{\frac{13}{25}}\,a}_{\red{r_{\text{质心}}}}\qquad T=\sqrt{\frac{4\pi^{2}a^{3}}{5Gm}} \] % 根式下方有红笔标注“r质心”；T 式分母的质量符号形似 m/M
%%%END
%%%BLOCK id=135-06 ch=99 sec=化学·零星 kind=note alt=none cont=none y=0.02-0.07
% 页眉右上角“Date.”处有一行淡灰色字迹（疑为相邻页内容；p133 顶部同一位置也有同样一行），与物理无关
\unclear{混合气体} $\rho=1.429\,\unit{g\cdot L^{-1}}$\quad（右端另有一处字迹 \unclear{$-6.72$}）
%%%END
%%%PAGEEND 135
